\documentclass[letterpaper,12pt,oneside]{article}
\usepackage[top=1in, left=1.25in, right=1.25in, bottom=1in]{geometry}

\usepackage[T1]{fontenc}
\usepackage[utf8]{inputenc}
\usepackage[spanish,es-nodecimaldot,es-tabla]{babel}
\usepackage{caption, subcaption} %figuras
\usepackage{graphicx}
\usepackage{array}
\usepackage{tikz}
\usepackage{imakeidx}
\usepackage{biblatex}
\addbibresource{bib/protocolo.bib}

\graphicspath{{./figs/}}
\usepackage{setspace}
%\usepackage[round]{natbib}

\begin{document}
%\frontmatter
\begin{titlepage}
		\setlength{\parindent}{0pt} \setlength{\parskip}{0pt}
	
		\begin{center}
			\vfill 
			
			\begin{minipage}{\textwidth}
				
				\newcolumntype{V}{>{\centering\arraybackslash} m{.17\textwidth} }
				\newcolumntype{C}{>{\centering\arraybackslash} m{.56\textwidth} }
				
				\begin{center}
					\includegraphics[width=4.5cm]{UNAM_LOGO2.png}	
				\end{center} 	
				\begin{center}
					\Huge\textbf{Universidad Nacional Autónoma de México}
				\end{center} 
				
			\end{minipage}
				
			\vfill
			
			\begin{minipage}{0.7\textwidth}
				\begin{center}
					\large Ingeniería en Computación
				\end{center}
			\end{minipage}
			% 
			\begin{center}
				%\medskip \rule{.9\textwidth}{2pt}
				\vfill
				\large Compiladores
				%{\Large \bfseries  \title   }
				
				%\medskip \rule{.9\textwidth}{2pt}
			\end{center}
			
			\begin{center}
				\vspace*{1cm}
				{\huge Analizador Léxico}\\
				\vspace*{1cm}
				
				%PRESENTA:\\ \medskip
				\begin{center}
					ALUMNO: \\
					\large 321184995\\
                    \large 000000000\\
                    \large 000000000\\
                    \large 000000000\\
                    \large 000000000
				\end{center}
				
				\bigskip
				
				Grupo: \\
				\#05\\
				%Coadvisors
				Semestre: \\
                2026-II
			\end{center}
			
			\vfill
			
			\begin{center}
				{México,CDMX. 15 Marzo 2025}\\
			\end{center}
			\cleardoublepage
		\end{center}
	\end{titlepage}
 
\tableofcontents
\clearpage
    
%\mainmatter

\section{Introducción} %

Este apartado debe abordar los siguientes puntos:

\begin{itemize}
    \item \textbf{Planteamiento del problema.} Se hace una descripción del problema a resolver.
    \item \textbf{Motivación.} Se describe porqué es necesario dar solución al problema.
    \item \textbf{Objetivos.} Lo que se se espera obtener de darle solución al problema.
\end{itemize}

\section{Marco Teórico}

Son los conceptos aplicados, principalmente los vistos en clase, con los que se va a dar solución al problema, explicados de forma concreta, ya sea con su interpretación personal o bien con referencias bibliográficas sólidas.~\cite{10.5555/576122}

\section{Desarrollo}

Es la descripción de la implementación realizada en el lenguaje de programación, así como las pruebas realizadas para obtener los resultados. Es importante resaltar lo siguiente:

\begin{itemize}
    \item No se debe mostrar código, solo describir funciones o la aplicación de los conceptos teóricos.
    \item Las pruebas es describir las entradas ingresadas y las salidas obtenidas.
\end{itemize}

\section{Resultados}

Mediante capturas de pantalla y una breve descripción seguida de la captura se presentan los resultados finales de su aplicación.

\section{Conclusiones}

Se presenta un análisis de los resultados obtenidos, donde se destaca la importancia de la aplicación de los conceptos teóricos para resolver el problema. Es importante resaltar lo siguiente:

\begin{itemize}
    \item No es describir si les gustó la actividad o no.
    \item No es decir que se obtuvo de la práctica.
    \item No es describir lo que fue difícil.
\end{itemize}

\printbibliography
\end{document}
